\documentclass[11pt]{article}

% ── Packages ──────────────────────────────────────────────────────────────────
\usepackage[margin=1in]{geometry}
\usepackage{amsmath,amssymb,amsthm}
\usepackage{booktabs}
\usepackage{algorithm}
\usepackage{algpseudocode}
\usepackage{graphicx}
\usepackage{hyperref}
\usepackage{xcolor}
\usepackage{multirow}
\usepackage{microtype}
% \usepackage{natbib}  % bibliography to be added by author

% ── Macros ────────────────────────────────────────────────────────────────────
\newcommand{\method}{EMA-GRPO}
\newcommand{\methodfull}{EMA-GRPO with Mastery Filter}
\newcommand{\pmb}{\mathcal{M}}
\newcommand{\ie}{\textit{i.e.}}
\newcommand{\eg}{\textit{e.g.}}
\newcommand{\etal}{\textit{et~al.}}

\theoremstyle{definition}
\newtheorem{definition}{Definition}

% ──────────────────────────────────────────────────────────────────────────────
\title{\textbf{Mitigating Catastrophic Forgetting in GRPO via\\
       Cross-Step EMA Normalization and Mastery Filtering}}

\author{
  Anonymous Author(s)
}

\date{}

\begin{document}

\maketitle

\begin{abstract}
Group Relative Policy Optimization (GRPO) has become a widely adopted
reinforcement learning algorithm for training large language models on
mathematical reasoning tasks.
However, we observe a systematic \emph{catastrophic forgetting} phenomenon:
questions answered correctly early in training regress to incorrect answers
in later epochs.
We identify three structural deficiencies in standard GRPO that collectively
cause this regression: (1) high-variance batch-level advantage normalization,
(2) destruction of the zero-gradient protection for mastered questions when
cross-step baselines are naively introduced, and (3) insufficient gradient
signal on hard problems under uniform sampling.
To address these deficiencies, we propose \methodfull{}, which maintains a
lightweight \emph{PromptMemoryBank} that accumulates per-prompt exponential
moving average (EMA) statistics across training steps.
Following \citet{liu2025drgrpo}, we do not divide by a standard deviation
estimate; instead the advantage is simply $r - \mu_{\mathrm{EMA}}$, which
reduces baseline variance by approximately $10\times$ compared to single-batch
statistics while avoiding std initialization artifacts.
A \emph{Mastery Filter} explicitly restores the zero-gradient protection by
zeroing the advantage of any prompt for which all current rollouts are correct,
preventing wasteful and harmful gradient updates on already-mastered content.
Experiments on MATH-500 and GSM8K with Qwen3-1.7B demonstrate consistent
improvements: +2.87\% pass@1 on MATH-500 and +0.79\% pass@1 on GSM8K,
with the forgetting collapse observed in standard GRPO entirely eliminated
on GSM8K.
\end{abstract}

% ══════════════════════════════════════════════════════════════════════════════
\section{Introduction}
\label{sec:intro}
% ══════════════════════════════════════════════════════════════════════════════

Reinforcement learning from verifiable rewards (RLVR) has emerged as a
powerful paradigm for eliciting mathematical reasoning from large language
models (LLMs)~\citep{ouyang2022training,lightman2023let}.
Among the algorithms proposed for this setting, Group Relative Policy
Optimization (GRPO)~\citep{shao2024deepseekmath} has gained broad adoption
due to its simplicity and effectiveness: it forgoes a learned critic by
normalizing rewards \emph{within} a group of rollouts sampled for the same
prompt, obtaining advantages without a separate value network.

Despite these advantages, we identify a systematic failure mode in GRPO that
has received limited attention: \textbf{catastrophic forgetting of mastered
problems during extended training}.
In our experiments with Qwen3-1.7B on MATH-500, 16 out of 373 training
problems that were correctly solved at epoch 1 regressed to incorrect answers
by epoch 9 under standard GRPO.
The regression patterns are qualitatively diverse: a model may correctly
derive intermediate steps but misidentify the target quantity at the final
step (\emph{last-step drift}), fall into repetitive verification loops that
diverge from the correct reasoning path (\emph{verification drift}), or
progressively drop boundary conditions that were present in earlier correct
responses (\emph{constraint erosion}).

We trace these forgetting phenomena to three interacting structural
deficiencies of standard GRPO:

\paragraph{Deficiency 1: High-variance batch-level advantage normalization.}
Standard GRPO normalizes advantages using the mean and standard deviation
computed over $n$ rollouts of the current batch.
With $n=8$ rollouts and a prompt whose true success probability is $\pi=0.5$,
the standard error of the batch mean is $\sqrt{\pi(1-\pi)/n} \approx 0.177$.
This noise causes the advantage baseline for the \emph{same} prompt to shift
substantially across training steps, producing contradictory gradient signals
that induce policy oscillation and forgetting.

\paragraph{Deficiency 2: Fragile zero-gradient protection.}
GRPO naturally produces zero gradient when all $n$ rollouts for a prompt are
correct (since all advantages are zero after within-group normalization),
providing implicit protection against over-training on mastered content.
Prior work has proposed cross-step baselines (\eg, EMA of historical rewards)
to stabilize normalization, but doing so naively breaks this protection:
an EMA mean of $\hat{\mu}=0.7$ assigns positive advantages to a fully-correct
rollout batch, causing the model to receive non-zero—and potentially
detrimental—gradient updates on already-mastered problems.
Our ablation confirms this: using EMA normalization alone \emph{increases}
the number of regressed problems from 19 to 37.

\paragraph{Deficiency 3: Insufficient gradient signal for hard problems.}
Under uniform prompt sampling, easy problems (high pass rate) and hard
problems (low pass rate) receive equal training exposure.
Because the per-step gradient signal for a problem scales with its advantage
magnitude, hard problems that the model rarely solves contribute relatively
weak learning signal, limiting the overall convergence rate.
As we show below, the Mastery Filter already addresses this deficiency
implicitly: by zeroing gradients on fully-mastered prompts, it ensures
that \emph{all} active gradient signal originates from problems the model
has not yet mastered---an implicit difficulty-aware curriculum
achieved without any explicit change to the sampling distribution.

To address all three deficiencies within a unified framework, we propose
\methodfull{}.
The core data structure is a \emph{PromptMemoryBank}: a lightweight dictionary
indexed by prompt ID that stores the per-prompt EMA mean reward and visit
count across training steps.
On top of this persistent memory, we build two complementary mechanisms:

\begin{enumerate}
  \item \textbf{EMA Mean Baseline.} We replace the noisy batch-level mean
    with a per-prompt EMA mean from the PromptMemoryBank, yielding a stable
    cross-step baseline: $A_i^{(p)} = r_i^{(p)} - \mu_{t-1}(p)$.
    Following \citet{liu2025drgrpo}, we do \emph{not} divide by a standard
    deviation estimate; the $n=8$ sample std is itself high-variance and
    its removal is consistently beneficial.
    With decay $\alpha=0.9$, the effective sample count of the mean grows to
    $n_{\text{eff}} \approx n/(1-\alpha) = 80$, reducing baseline variance
    by $\approx10\times$ relative to standard GRPO.

  \item \textbf{Mastery Filter.} We explicitly restore the zero-gradient
    protection by hard-zeroing the advantage of any prompt for which the
    current rollout batch is fully correct (\ie, all $n$ rollouts succeed).
    Crucially, the PromptMemoryBank still records the correct-batch mean
    for future EMA updates, so the memory stays current even when no gradient
    is emitted.
    Beyond protecting mastered knowledge, the Filter has a secondary effect:
    it concentrates \emph{all} gradient signal on the subset of prompts that
    the model has not yet mastered, forming an implicit difficulty-aware
    curriculum that also addresses Deficiency~3 without any explicit
    change to the sampling distribution.
\end{enumerate}

We explored an additional \textbf{Degradation-Aware Resampling} component
that explicitly upweights hard and regressing prompts at the sampling level.
However, this is conceptually redundant with the Mastery Filter's implicit
gradient concentration: the Filter already ensures that every gradient step
draws exclusively from non-mastered prompts.
Our ablation (Section~\ref{sec:resampling_ablation}) confirms that explicit
resampling provides no additional benefit and is harmful on smaller training
sets due to overfitting amplification.

We evaluate \method{} on two benchmarks.
On MATH-500 (400 training problems), \method{} achieves pass@1 of 0.689 versus
baseline GRPO's 0.670 (+2.87\%), and pass@8 of 0.800 versus 0.790 (+1.27\%).
On GSM8K (1{,}000 training problems), standard GRPO reaches a peak of 92.16\%
accuracy before collapsing to 88.09\% by the end of training—a canonical
forgetting trajectory. \method{} instead climbs steadily to 92.48\% and
remains stable at 89.34\% at the same final step, with no collapse.

Our contributions are as follows.

\paragraph{(1) Forgetting characterization.}
We provide a detailed empirical characterization of catastrophic forgetting
in GRPO, identifying three interacting structural root causes and documenting
qualitative regression patterns observed in extended training.

\paragraph{(2) PromptMemoryBank.}
We introduce the \emph{PromptMemoryBank} abstraction: a lightweight
per-prompt EMA statistics store that persists across training steps,
enabling cross-step advantage normalization without breaking
zero-gradient protection.

\paragraph{(3) \methodfull{} algorithm.}
We combine the EMA mean baseline and the Mastery Filter into a principled
algorithm, supported by thorough component ablations that validate each design
choice and rule out several intuitive but harmful alternatives.

\paragraph{(4) Empirical validation.}
We demonstrate consistent forgetting mitigation and accuracy gains across
four evaluation benchmarks (AIME 2024, AIME 2025, MATH-500, GSM8K),
two model families, and three training set sizes.

% ══════════════════════════════════════════════════════════════════════════════
\section{Related Work}
\label{sec:related}
% ══════════════════════════════════════════════════════════════════════════════

\subsection{Policy Optimization for LLM Reasoning}

\paragraph{PPO and RLHF.}
Proximal Policy Optimization (PPO)~\citep{schulman2017proximal} is the
canonical on-policy RL algorithm for LLM alignment~\citep{ouyang2022training}.
PPO relies on a learned value network to compute advantages, which adds
significant memory and compute overhead in the LLM setting.
REINFORCE++~\citep{hu2025reinforce} revisits the simpler REINFORCE estimator
with a token-level baseline to reduce variance, showing competitive performance
without a critic network.

\paragraph{GRPO.}
\citet{shao2024deepseekmath} proposed GRPO as a critic-free alternative:
for each prompt, $n$ responses are sampled, and the reward of each response
is normalized by the mean and standard deviation of the group to obtain
advantages.
This eliminates the value network while providing relative reward signals.
GRPO has since been adopted by a wide range of mathematical reasoning
systems~\citep{deepseekr1,qwen2math}.

\paragraph{DAPO.}
\citet{yu2025dapo} identified that all-correct and all-wrong rollout batches
produce zero gradient in GRPO, and proposed \emph{Dynamic Sampling} to
discard such batches and resample from the remaining prompts.
While this eliminates zero-gradient steps, it discards all information from
mastered prompts and does not accumulate any cross-step statistics.
Our Mastery Filter takes a different position: instead of discarding
all-correct batches, we zero their gradient contribution while \emph{retaining}
them for PromptMemoryBank updates.
This preserves the EMA statistics required for stable future normalization.

\paragraph{Dr.\ GRPO.}
\citet{liu2025drgrpo} observed that batch-level normalization in GRPO
conflates per-prompt difficulty, introducing systematic bias.
They proposed removing the std normalization entirely, computing advantages
as $A_i^{(p)} = r_i^{(p)} - \hat{\mu}_t^{(p)}$ using only the per-prompt
batch mean as baseline.
This eliminates the noise from the $n=8$ sample std estimator.
Our method shares the no-std design principle with Dr.\ GRPO, but replaces
the single-step batch mean with a cross-step EMA mean: the EMA baseline
is an order of magnitude less noisy ($n_{\mathrm{eff}} = 80$ vs.\ $n=8$)
and explicitly tracks per-prompt performance history, enabling the Mastery
Filter and Degradation-Aware Resampling mechanisms.

\subsection{Catastrophic Forgetting in Neural Networks}

Catastrophic forgetting, the tendency of neural networks to lose previously
acquired knowledge upon learning new information, is a classical problem in
continual learning~\citep{mccloskey1989catastrophic,french1999catastrophic}.
Standard mitigations include Elastic Weight Consolidation
(EWC)~\citep{kirkpatrick2017overcoming}, which penalizes changes to weights
important for past tasks; Progressive Neural Networks~\citep{rusu2016progressive},
which allocate new capacity for new tasks; and experience
replay~\citep{riemer2018learning,rolnick2019experience}, which maintains a
buffer of past examples to interleave with new training.

In the LLM fine-tuning context, forgetting of pre-training capabilities
during instruction tuning has been studied by
\citet{luo2023empirical,kotha2023understanding}.
KL-divergence penalties against a frozen reference model—as used in standard
RLHF and in GRPO's regularization term—provide a global constraint on policy
drift but cannot selectively protect specific capabilities.
Our Mastery Filter offers a different mechanism: rather than constraining
\emph{how much} the policy can change globally, it identifies specific mastered
prompts and suppresses gradient updates on them precisely.

\subsection{Curriculum Learning and Adaptive Sampling}

Curriculum learning~\citep{bengio2009curriculum} proposes to order training
examples from easy to hard, accelerating learning and sometimes improving
final performance.
Self-paced learning~\citep{kumar2010self} and its variants automate curriculum
construction by weighting examples according to current model difficulty.
In the RL-for-reasoning literature, \citet{fan2025areal} and
\citet{zhao2025absolute} explore difficulty-based sampling strategies.
Our Degradation-Aware Resampling combines difficulty-based weighting with
regression detection: a prompt's sampling weight increases both when it is
historically hard \emph{and} when it has recently regressed, targeting
precisely the problems that represent active forgetting events.

\subsection{EMA in Optimization and RL}

Exponential moving averages are ubiquitous in optimization, from Adam's
moment estimates~\citep{kingma2014adam} to target networks in
DQN~\citep{mnih2015human} and soft actor-critic~\citep{haarnoja2018soft}.
In the context of advantage normalization, \citet{andrychowicz2020matters}
discuss running statistics for reward normalization in deep RL.
Our use of EMA is most closely related to running reward normalization, but
differs in operating at the \emph{per-prompt} level rather than globally,
and in being coupled with the Mastery Filter to prevent the zero-gradient
protection failure described in Section~\ref{sec:intro}.

% ══════════════════════════════════════════════════════════════════════════════
\section{Method}
\label{sec:method}
% ══════════════════════════════════════════════════════════════════════════════

\subsection{Preliminaries: GRPO}
\label{sec:grpo_prelim}

Let $\mathcal{D} = \{p_1, \ldots, p_N\}$ be a training set of prompts.
At each training step $t$, a mini-batch $\mathcal{B} \subset \mathcal{D}$ is
sampled, and for each prompt $p \in \mathcal{B}$, $n$ responses are rolled out
from the current policy $\pi_\theta$:
$\{o_1^{(p)}, \ldots, o_n^{(p)}\}$.
Each response receives a scalar reward $r_i^{(p)} \in [0,1]$ from a verifiable
reward function (e.g., exact-match against ground-truth answer).
Standard GRPO computes the advantage of response $o_i^{(p)}$ as
\begin{equation}
  A_i^{(p)} = \frac{r_i^{(p)} - \hat{\mu}_t^{(p)}}{\hat{\sigma}_t^{(p)} + \varepsilon},
  \label{eq:grpo_adv}
\end{equation}
where $\hat{\mu}_t^{(p)} = \frac{1}{n}\sum_i r_i^{(p)}$ and
$\hat{\sigma}_t^{(p)} = \mathrm{std}(\{r_i^{(p)}\})$ are \emph{batch-level}
statistics computed from the current step's $n$ rollouts only.
The policy is updated by maximizing the clipped surrogate objective
\begin{equation}
  \mathcal{L}_{\text{GRPO}}(\theta) =
  \mathbb{E}\!\left[\min\!\left(
    \frac{\pi_\theta(o|p)}{\pi_{\theta_{\text{old}}}(o|p)} A,\;
    \mathrm{clip}\!\left(\frac{\pi_\theta(o|p)}{\pi_{\theta_{\text{old}}}(o|p)},
      1\!-\!\epsilon, 1\!+\!\epsilon\right) A
  \right)\right]
  - \beta\, D_{\mathrm{KL}}(\pi_\theta \| \pi_{\mathrm{ref}}),
  \label{eq:grpo_obj}
\end{equation}
with clip ratio $\epsilon$ and KL coefficient $\beta$.

\subsection{PromptMemoryBank}
\label{sec:pmb}

The central data structure of our method is the \emph{PromptMemoryBank}
$\mathcal{M}$, a dictionary keyed by prompt index $\mathrm{id}(p) \in \{1,\ldots,N\}$.
Each entry stores a pair
$\mathcal{M}[p] = (\mu(p),\, c(p))$,
where $\mu(p)$ is the EMA-smoothed mean reward for prompt $p$ and $c(p)$
is the visit count.
We intentionally do not track a per-prompt EMA standard deviation:
dividing by the std estimate introduces initialization artifacts (the
standard prior of $\sigma_0=1$ takes $\sim\!20$ visits to converge to the
true value of $\approx\!0.5$ for binary rewards) and, per \citet{liu2025drgrpo},
is not necessary for effective advantage estimation with small rollout groups.

\paragraph{EMA update rule.}
After computing the batch mean $\hat{\mu}_t^{(p)}$ from the current step's
$n$ rollouts, we update $\mathcal{M}[p]$ using
\begin{equation}
  \mu_t(p) \leftarrow \alpha\,\mu_{t-1}(p) + (1-\alpha)\,\hat{\mu}_t^{(p)},
  \label{eq:ema_mean}
\end{equation}
where $\alpha \in (0,1)$ is the EMA decay coefficient (default $\alpha=0.9$).

\paragraph{Cold-start initialization.}
On the \emph{first} visit to prompt $p$ ($c(p)=0$), we initialize
$\mu_0(p) \leftarrow \hat{\mu}_t^{(p)}$ directly (cold start) rather than
using the arbitrary prior $\mu_0(p)=0$.
This ensures the EMA baseline is grounded in actual model performance from
the very first encounter, eliminating the warm-up period of mean underestimation.

\paragraph{Effective sample count.}
The EMA with decay $\alpha$ is equivalent to a weighted average over past
batches with exponentially decaying weights, yielding an effective sample
count of
\begin{equation}
  n_{\mathrm{eff}} \approx \frac{n}{1-\alpha}.
\end{equation}
With $n=8$ and $\alpha=0.9$, $n_{\mathrm{eff}} = 80$, reducing the variance
of the mean estimator by a factor of $10$ relative to standard GRPO.

\paragraph{Warmup.}
Although the cold-start initialization grounds the EMA mean from the first
visit, the estimate can still be noisy after a single observation.
We introduce a warmup threshold $K$: for visits $c(p) < K$, advantage
computation falls back to the current batch mean as in standard GRPO.
EMA normalization activates only after $c(p) \geq K$ (default $K=3$).

\subsection{EMA Mean Baseline}
\label{sec:ema_norm}

Once the warmup condition is satisfied, we subtract the EMA mean from the
PromptMemoryBank as the advantage baseline, without any standard deviation
scaling:
\begin{equation}
  A_i^{(p)} = r_i^{(p)} - \mu_{t-1}(p),
  \qquad c(p) \geq K.
  \label{eq:ema_adv}
\end{equation}
We use $\mu_{t-1}(p)$ (\ie, the EMA mean \emph{before} the current-step
update) to avoid look-ahead bias: the advantage is computed before the memory
bank incorporates the new batch.
The EMA update (Eq.~\eqref{eq:ema_mean}) is applied \emph{after} the
advantage is computed.

The omission of std normalization is deliberate and follows from three
complementary arguments.

\textbf{(i) For binary rewards, std carries no information beyond the mean.}
Mathematical reasoning tasks use binary correctness rewards $r \in \{0,1\}$.
For a Bernoulli random variable with mean $\mu$, the standard deviation is
exactly $\sigma = \sqrt{\mu(1-\mu)}$---a deterministic function of $\mu$.
Given the EMA mean $\mu_{t-1}(p)$, the per-prompt std is therefore fully
determined: dividing by $\hat{\sigma}$ merely applies a fixed monotone
rescaling $1/\sqrt{\mu(1-\mu)}$ to the advantage, without adding any
information about the prompt's difficulty or the quality of the current
response.
Concretely, Eq.~\eqref{eq:ema_adv} gives $A_i^{(p)} \in \{1-\mu, -\mu\}$
for a correct and incorrect response respectively---magnitudes that naturally
reflect difficulty: a hard prompt ($\mu \approx 0$) yields a large positive
advantage when solved ($\approx 1$) and a small negative when failed
($\approx 0$), concentrating gradient energy where learning is most needed.
Dividing by $\sigma$ would distort this natural weighting without benefit.

\textbf{(ii) Small-sample std estimates are high-variance.}
With $n=8$ rollouts, the empirical std $\hat{\sigma}_t^{(p)}$ is an
unreliable estimate of the true $\sigma(p)$.
Its coefficient of variation is $\approx 1/\!\sqrt{2(n-1)} \approx 0.27$
for $n=8$---meaning the denominator can swing by $\pm 27\%$ step-to-step
from sampling noise alone.
Dividing by this estimate re-introduces step-level variance into the
advantage after the EMA mean has removed it.
\citet{liu2025drgrpo} empirically confirms that removing std normalization
is consistently beneficial for small rollout groups.
An EMA std would not help, and in fact introduces a critical instability.
Consider the prior $\sigma_0 = 1$: this far exceeds the true binary-reward
std of $\approx 0.5$, and with $\alpha=0.9$ requires $\sim\!20$ EMA
updates to converge---producing an abrupt advantage-scale shift when EMA
activates that destabilises training.
More dangerously, if a prompt's first $n$ rollouts are \emph{all correct}
or \emph{all incorrect} ($\hat{\sigma}=0$), a cold-start initialization
of the EMA std stores $\hat{\sigma}=\epsilon \approx 10^{-6}$ to avoid
division by zero.
The EMA decay preserves this near-zero value: after $k$ subsequent updates,
$\sigma_k \approx \alpha^k \epsilon + (1-\alpha^k)\,\bar{\sigma}$.
With $\alpha=0.9$ and $\epsilon=10^{-6}$, even after 10 updates
$\sigma_{10} \approx 0.35\,\epsilon + 0.65\,\bar{\sigma} \approx 0.32$---still
well below the batch std.
When the warmup ends and EMA activates, advantages are divided by this
underestimated $\sigma$, causing a sudden $\sim\!2\times$ scale increase.
In practice we observed $\mathtt{pg\_loss}$ spikes to values in the
range $10^3$--$10^4$ at the warmup boundary, followed by gradient explosion
that destabilises or terminates training.
Removing std from the EMA entirely eliminates this failure mode at its root.

\textbf{(iii) The Mastery Filter handles the degenerate $\sigma \to 0$ case.}
The canonical reason to keep a std in the denominator is numerical safety:
when $\sigma \approx 0$ (the model is nearly deterministic on a prompt),
division could explode.
The Mastery Filter eliminates this concern: it zeroes the advantage for
any prompt where all $n$ rollouts are correct ($\hat{\mu}_t^{(p)} = 1.0$,
$\hat{\sigma}_t^{(p)} = 0$), the only case where $\sigma = 0$ for binary
rewards.
The two mechanisms are therefore complementary: the Mastery Filter makes
std division unnecessary for numerical safety, and the binary-reward
structure makes it unnecessary for information content.

\subsection{Mastery Filter}
\label{sec:mastery}

As noted in Section~\ref{sec:intro}, EMA normalization alone breaks the
zero-gradient protection of standard GRPO.
When $\mu_{t-1}(p)$ lags behind the true current performance (\eg,
$\mu_{t-1}(p) = 0.7$ while all $n$ current rollouts succeed), the advantages
in Eq.~\eqref{eq:ema_adv} are positive, causing the policy to receive a
non-zero update that pushes it \emph{further} toward responses it already
produces reliably—a form of gradient waste that can disrupt neighboring
capabilities.

The Mastery Filter restores exact zero-gradient protection by detecting
fully-correct batches and zeroing their advantages:
\begin{equation}
  A_i^{(p)} = 0 \quad \text{if } \hat{\mu}_t^{(p)} \geq \tau_{\mathrm{mastery}},
  \label{eq:mastery}
\end{equation}
where $\tau_{\mathrm{mastery}} = 1.0$ (all $n$ rollouts correct).
We set $\tau = 1.0$ rather than a softer threshold because a $(n-1)/n$-correct
batch still contains one failed rollout that provides a valid contrastive
gradient signal.

Crucially, zeroing the advantage does \emph{not} skip the PromptMemoryBank
update: the EMA mean $\mu(p)$ is updated regardless, staying current even
during mastered-prompt steps.
This ensures that when the model later regresses on a previously mastered
problem, the EMA mean quickly reflects the regression and the Mastery
Filter disengages, allowing corrective gradient updates to resume.

\paragraph{Implicit difficulty-aware curriculum.}
The Mastery Filter has a second, less obvious effect: it acts as an
implicit difficulty-aware curriculum learning mechanism in gradient space.
Because all fully-correct prompts contribute zero gradient, \emph{every
non-zero gradient step is driven entirely by prompts the model has not yet
mastered}---hard problems (consistently low $\mu$) and recently regressed
problems (falling $\mu$) both retain full gradient weight, while easy
mastered problems are silenced.
This implicit curriculum emerges automatically from the filter's primary
forgetting-prevention role, without any explicit change to the training
data sampling distribution.
An explicit Degradation-Aware Resampling scheme (Section~\ref{sec:resampling})
targets the same objective at the sampling level, but our experiments show
it provides no additional gain over the implicit curriculum already established
by the Mastery Filter, and causes overfitting on small datasets.

\paragraph{Combined advantage computation.}
Putting everything together, the advantage for response $o_i^{(p)}$ at step
$t$ is:
\begin{equation}
  A_i^{(p)} = \begin{cases}
    0
      & \text{if } \hat{\mu}_t^{(p)} \geq \tau_{\mathrm{mastery}}
        \quad\text{(Mastery Filter)}\\[4pt]
    r_i^{(p)} - \mu_{t-1}(p)
      & \text{if } c(p) \geq K
        \quad\text{(EMA mean baseline)}\\[4pt]
    r_i^{(p)} - \hat{\mu}_t^{(p)}
      & \text{otherwise}
        \quad\text{(warmup: batch mean baseline)}
  \end{cases}
  \label{eq:final_adv}
\end{equation}

\subsection{Full Algorithm}
\label{sec:algorithm}

Algorithm~\ref{alg:ema_grpo} describes the complete \method{} training loop.

\begin{algorithm}[t]
\caption{\methodfull}
\label{alg:ema_grpo}
\begin{algorithmic}[1]
\Require Training set $\mathcal{D}$, policy $\pi_\theta$, reference $\pi_{\mathrm{ref}}$,
         EMA decay $\alpha$, warmup $K$, mastery threshold $\tau$, rollout count $n$
\State Initialize $\mathcal{M}[p] \leftarrow \textit{empty}$ for all $p \in \mathcal{D}$
\For{each training step $t$}
  \State Sample mini-batch $\mathcal{B} \subset \mathcal{D}$
  \For{each prompt $p \in \mathcal{B}$}
    \State Roll out $n$ responses; compute rewards $\{r_i^{(p)}\}$
    \State $\hat{\mu}_t^{(p)} \leftarrow \frac{1}{n}\sum_i r_i^{(p)}$
    \State Read $(\mu_{t-1}(p),\; c(p)) \leftarrow \mathcal{M}[p]$
    \For{$i = 1, \ldots, n$}
      \If{$\hat{\mu}_t^{(p)} \geq \tau$} \hfill\Comment{Mastery Filter}
        \State $A_i^{(p)} \leftarrow 0$
      \ElsIf{$c(p) \geq K$} \hfill\Comment{EMA baseline}
        \State $A_i^{(p)} \leftarrow r_i^{(p)} - \mu_{t-1}(p)$
      \Else \hfill\Comment{warmup}
        \State $A_i^{(p)} \leftarrow r_i^{(p)} - \hat{\mu}_t^{(p)}$
      \EndIf
    \EndFor
    \State Update $\mathcal{M}[p]$ via Eq.~\eqref{eq:ema_mean}; increment $c(p)$
  \EndFor
  \State Update $\pi_\theta$ with objective~\eqref{eq:grpo_obj} using $\{A_i^{(p)}\}$
\EndFor
\end{algorithmic}
\end{algorithm}

% ══════════════════════════════════════════════════════════════════════════════
\section{Experiments}
\label{sec:experiments}
% ══════════════════════════════════════════════════════════════════════════════

\subsection{Experimental Setup}
\label{sec:setup}

\paragraph{Models.}
We evaluate across four model configurations of increasing capacity:
\textbf{Qwen3-1.7B} (used for ablation studies),
\textbf{Qwen2.5-Math-1.5B} (primary main results),
\textbf{Qwen3-8B} (additional validation), and
\textbf{Qwen2.5-Math-7B-Base} (large-scale generalization, Section~\ref{sec:largescale}).
All models are fine-tuned with LoRA (rank~64, $\alpha_{\mathrm{LoRA}}=32$).

\paragraph{Training data.}
We train on \textbf{MATH-Train$_{\text{2800}}$}, a 2{,}800-problem subset
sampled from the MATH training split.
Smaller subsets used in ablation studies (400-problem and 1{,}000-problem splits)
are described in the respective subsections.

\paragraph{Evaluation benchmarks.}
All methods are evaluated on four benchmarks across the full difficulty spectrum:
\textbf{MATH-500} (500 problems, broad mathematical reasoning),
\textbf{GSM8K} (1{,}319 problems, grade-school arithmetic),
\textbf{AIME 2024} (30 competition problems, high difficulty), and
\textbf{AIME 2025} (30 competition problems, post-training cutoff, out-of-distribution).
For each benchmark we report pass@1, pass@4, and pass@8, estimated with
$n=8$ samples at temperature~0.6 using the unbiased estimator.
AIME 2025 post-dates all training data, providing a clean cross-distribution
generalization signal.

\paragraph{Hyperparameters.}
All models use LoRA (rank~64, $\alpha_{\mathrm{LoRA}}=32$), AdamW with
lr~$=10^{-5}$, batch size~32, rollout count $n=8$, and max response
length 8{,}192 tokens.
The KL coefficient is $\beta=0.001$ (low-variance estimator).
For \method{} we use EMA decay $\alpha=0.9$, warmup threshold $K=3$,
and mastery threshold $\tau=1.0$; ablations for these choices are in
Section~\ref{sec:hparam} and Appendix~\ref{app:impl}.
Full hyperparameter listings are provided in Table~\ref{tab:hyperparams_full}
(Appendix~\ref{app:impl}).

\paragraph{Baselines.}
We compare against three baselines.
\textbf{GRPO} is the standard baseline with batch-level mean and std normalization.
\textbf{Dr.~GRPO} removes std normalization while retaining the batch-level mean baseline.
\textbf{DAPO} adds dynamic group filtering (discards all-correct/all-wrong batches),
clip-higher ($\epsilon_{\mathrm{high}}=0.28$), entropy bonus ($0.001$),
and removes the KL penalty.

% ──────────────────────────────────────────────────────────────────────────────
\subsection{Main Results}
\label{sec:main_results}
% ──────────────────────────────────────────────────────────────────────────────

Tables~\ref{tab:main_final} and~\ref{tab:main_best} report results on all four
evaluation benchmarks for methods trained on MATH-Train$_{\text{2800}}$
with Qwen2.5-Math-1.5B and Qwen3-8B.
Table~\ref{tab:main_final} uses the \emph{final} checkpoint (step~1720, epoch~20);
Table~\ref{tab:main_best} uses the \emph{best intermediate} checkpoint
per method, selected by AIME~2024 pass@1, to isolate peak performance
from late-training regression.
Figure~\ref{fig:training_curves} visualizes training dynamics across all benchmarks.

\begin{table}[t]
\centering
\caption{Main results at the \textbf{final} checkpoint (step~1720, epoch~20)
  on MATH-Train$_{\text{2800}}$.
  Best per column within each model group in \textbf{bold}.}
\label{tab:main_final}
\resizebox{\linewidth}{!}{%
\begin{tabular}{llcccccccccccc}
\toprule
& & \multicolumn{3}{c}{\textbf{AIME 2024}} & \multicolumn{3}{c}{\textbf{AIME 2025}}
& \multicolumn{3}{c}{\textbf{MATH-500}} & \multicolumn{3}{c}{\textbf{GSM8K}} \\
\cmidrule(lr){3-5}\cmidrule(lr){6-8}\cmidrule(lr){9-11}\cmidrule(lr){12-14}
\textbf{Model} & \textbf{Method} & p@1 & p@4 & p@8 & p@1 & p@4 & p@8 & p@1 & p@4 & p@8 & p@1 & p@4 & p@8 \\
\midrule
\multirow{4}{*}{Qwen2.5-Math-1.5B}
  & GRPO             & -- & -- & -- & -- & -- & -- & -- & -- & -- & -- & -- & -- \\
  & Dr.~GRPO         & -- & -- & -- & -- & -- & -- & -- & -- & -- & -- & -- & -- \\
  & DAPO             & -- & -- & -- & -- & -- & -- & -- & -- & -- & -- & -- & -- \\
  & \method{} (ours) & -- & -- & -- & -- & -- & -- & -- & -- & -- & -- & -- & -- \\
\midrule
\multirow{4}{*}{Qwen3-8B}
  & GRPO             & -- & -- & -- & -- & -- & -- & -- & -- & -- & -- & -- & -- \\
  & Dr.~GRPO         & -- & -- & -- & -- & -- & -- & -- & -- & -- & -- & -- & -- \\
  & DAPO             & -- & -- & -- & -- & -- & -- & -- & -- & -- & -- & -- & -- \\
  & \method{} (ours) & -- & -- & -- & -- & -- & -- & -- & -- & -- & -- & -- & -- \\
\bottomrule
\end{tabular}}
\end{table}

\begin{table}[t]
\centering
\caption{Main results at the \textbf{best intermediate} checkpoint
  (selected by AIME~2024 pass@1; step in parentheses).
  Same training setup as Table~\ref{tab:main_final}.}
\label{tab:main_best}
\resizebox{\linewidth}{!}{%
\begin{tabular}{llccccccccccccc}
\toprule
& & & \multicolumn{3}{c}{\textbf{AIME 2024}} & \multicolumn{3}{c}{\textbf{AIME 2025}}
& \multicolumn{3}{c}{\textbf{MATH-500}} & \multicolumn{3}{c}{\textbf{GSM8K}} \\
\cmidrule(lr){4-6}\cmidrule(lr){7-9}\cmidrule(lr){10-12}\cmidrule(lr){13-15}
\textbf{Model} & \textbf{Method} & \textbf{Step} & p@1 & p@4 & p@8 & p@1 & p@4 & p@8 & p@1 & p@4 & p@8 & p@1 & p@4 & p@8 \\
\midrule
\multirow{4}{*}{Qwen2.5-Math-1.5B}
  & GRPO             & -- & -- & -- & -- & -- & -- & -- & -- & -- & -- & -- & -- & -- \\
  & Dr.~GRPO         & -- & -- & -- & -- & -- & -- & -- & -- & -- & -- & -- & -- & -- \\
  & DAPO             & -- & -- & -- & -- & -- & -- & -- & -- & -- & -- & -- & -- & -- \\
  & \method{} (ours) & -- & -- & -- & -- & -- & -- & -- & -- & -- & -- & -- & -- & -- \\
\midrule
\multirow{4}{*}{Qwen3-8B}
  & GRPO             & -- & -- & -- & -- & -- & -- & -- & -- & -- & -- & -- & -- & -- \\
  & Dr.~GRPO         & -- & -- & -- & -- & -- & -- & -- & -- & -- & -- & -- & -- & -- \\
  & DAPO             & -- & -- & -- & -- & -- & -- & -- & -- & -- & -- & -- & -- & -- \\
  & \method{} (ours) & -- & -- & -- & -- & -- & -- & -- & -- & -- & -- & -- & -- & -- \\
\bottomrule
\end{tabular}}
\end{table}

\begin{figure}[t]
\centering
\fbox{\parbox{0.95\linewidth}{\centering\vspace{3.5cm}
  \textit{[Placeholder: $2{\times}2$ grid of training accuracy curves.
  Subplots (left-to-right, top-to-bottom): AIME~2024 pass@1,
  AIME~2025 pass@1, MATH-500 pass@1, GSM8K pass@1.
  Each subplot shows curves for GRPO, Dr.~GRPO, DAPO, and \method{}
  over training steps 0--1720.
  Annotate the forgetting collapse in GRPO (sharp drop after peak)
  and the stable trajectory of \method{}.]}
  \vspace{3.5cm}}}
\caption{Training accuracy curves on all four evaluation benchmarks.
  \method{} maintains a stable improvement trajectory; standard GRPO
  exhibits a characteristic forgetting collapse after reaching peak accuracy.}
\label{fig:training_curves}
\end{figure}

% ──────────────────────────────────────────────────────────────────────────────
\subsection{Ablation Studies}
\label{sec:ablation}
% ──────────────────────────────────────────────────────────────────────────────

\subsubsection{Component Ablation}
\label{sec:comp_ablation}

To isolate the contribution of each design choice, we ablate three binary
switches: (1) replacing the batch mean with an EMA mean, (2) applying the
Mastery Filter, and (3) removing std normalization.
Experiments use Qwen3-1.7B trained on MATH-500$_{\text{400}}$.
Table~\ref{tab:ablation_component} reports best-checkpoint pass@1 and pass@8
for all four benchmarks.

A critical finding is that EMA normalization \emph{alone} (without the
Mastery Filter) increases the number of regressed training problems from
19 to 37, confirming the zero-gradient protection failure described in
Section~\ref{sec:intro}.
The Mastery Filter restores this protection and yields the strongest results.

\begin{table}[t]
\centering
\caption{Component ablation (Qwen3-1.7B, MATH-500$_{\text{400}}$, best checkpoint pass@1 / pass@8).}
\label{tab:ablation_component}
\resizebox{\linewidth}{!}{%
\begin{tabular}{lcccccccccc}
\toprule
& & & & \multicolumn{2}{c}{\textbf{AIME 2024}} & \multicolumn{2}{c}{\textbf{AIME 2025}}
& \multicolumn{2}{c}{\textbf{MATH-500}} & \multicolumn{2}{c}{\textbf{GSM8K}} \\
\cmidrule(lr){5-6}\cmidrule(lr){7-8}\cmidrule(lr){9-10}\cmidrule(lr){11-12}
\textbf{Variant} & \textbf{EMA} & \textbf{Filter} & \textbf{No std}
  & p@1 & p@8 & p@1 & p@8 & p@1 & p@8 & p@1 & p@8 \\
\midrule
GRPO (baseline)  & \texttimes & \texttimes & \texttimes & -- & -- & -- & -- & -- & -- & -- & -- \\
+~EMA only       & \checkmark & \texttimes & \texttimes & -- & -- & -- & -- & -- & -- & -- & -- \\
+~EMA + Filter   & \checkmark & \checkmark & \texttimes & -- & -- & -- & -- & -- & -- & -- & -- \\
\method{} (ours) & \checkmark & \checkmark & \checkmark & -- & -- & -- & -- & -- & -- & -- & -- \\
\bottomrule
\end{tabular}}
\end{table}

\subsubsection{Effect of Std Normalization}
\label{sec:std_ablation}

Table~\ref{tab:ablation_std} isolates the impact of dividing by a std
estimate, comparing four advantage formulations that vary only in the
source of the mean and std.
The ``EMA std (w/ bug)'' row corresponds to an earlier implementation that
tracked EMA std with a $\sigma_0=1$ prior, causing a transient
pg\_loss spike when EMA activates (as analyzed in Section~\ref{sec:pmb}).
This ablation, run on Qwen2.5-Math-1.5B with MATH-Train$_{\text{2800}}$,
tests whether removing std normalization is strictly beneficial in combination
with the EMA mean baseline.

\begin{table}[t]
\centering
\caption{Std normalization ablation (Qwen2.5-Math-1.5B, MATH-Train$_{\text{2800}}$,
  best checkpoint).
  ``Batch'' = current-step rollout statistics;
  ``EMA'' = cross-step moving average;
  ``--'' = not applied.}
\label{tab:ablation_std}
\resizebox{\linewidth}{!}{%
\begin{tabular}{llcccccccc}
\toprule
& & \multicolumn{2}{c}{\textbf{AIME 2024}} & \multicolumn{2}{c}{\textbf{AIME 2025}}
& \multicolumn{2}{c}{\textbf{MATH-500}} & \multicolumn{2}{c}{\textbf{GSM8K}} \\
\cmidrule(lr){3-4}\cmidrule(lr){5-6}\cmidrule(lr){7-8}\cmidrule(lr){9-10}
\textbf{Method} & \textbf{Mean / Std}
  & p@1 & p@8 & p@1 & p@8 & p@1 & p@8 & p@1 & p@8 \\
\midrule
GRPO                   & Batch / Batch & -- & -- & -- & -- & -- & -- & -- & -- \\
Dr.~GRPO               & Batch / --    & -- & -- & -- & -- & -- & -- & -- & -- \\
EMA-GRPO (w/ EMA std)  & EMA / EMA     & -- & -- & -- & -- & -- & -- & -- & -- \\
\method{} (ours)       & EMA / --      & -- & -- & -- & -- & -- & -- & -- & -- \\
\bottomrule
\end{tabular}}
\end{table}

\subsubsection{Forgetting Analysis}
\label{sec:forgetting}

To directly quantify catastrophic forgetting, we track per-prompt pass@1
across epochs and count prompts that regress from mastered ($\geq 0.5$
pass@1 at epoch~1) to forgotten ($< 0.5$ pass@1 at epoch~20).
Table~\ref{tab:forgetting} reports regression statistics for all methods.
Figures~\ref{fig:heatmap} and~\ref{fig:regression_curve} visualize
the forgetting dynamics in detail.

\begin{table}[t]
\centering
\caption{Forgetting quantification: prompts regressing from pass@1$\geq$0.5
  at epoch~1 to pass@1$<$0.5 at epoch~20.
  Qwen2.5-Math-1.5B, MATH-Train$_{\text{2800}}$ (2{,}800 prompts total).}
\label{tab:forgetting}
\small
\begin{tabular}{lccc}
\toprule
\textbf{Method} & \textbf{\# Regressed} & \textbf{\% of training set} & \textbf{Max regression depth} \\
\midrule
GRPO             & -- & -- & -- \\
Dr.~GRPO         & -- & -- & -- \\
DAPO             & -- & -- & -- \\
\method{} (ours) & -- & -- & -- \\
\bottomrule
\end{tabular}
\end{table}

\begin{figure}[t]
\centering
\fbox{\parbox{0.95\linewidth}{\centering\vspace{3.5cm}
  \textit{[Placeholder: side-by-side per-prompt forgetting heatmaps.
  Left: GRPO. Right: \method{}.
  X-axis: training step (0--1720, checkpointed every 200 steps).
  Y-axis: training prompts sorted by initial difficulty (hardest top).
  Color: pass@1 at each checkpoint (green = correct, red = incorrect).
  Highlight regression events (green $\to$ red) that appear in GRPO
  but are suppressed or recovered in \method{}.]}
  \vspace{3.5cm}}}
\caption{Per-prompt pass@1 heatmaps over training steps for GRPO (left)
  vs.\ \method{} (right). Each row is one training prompt.
  Regression events (green $\to$ red transitions) are visibly suppressed
  by the Mastery Filter.}
\label{fig:heatmap}
\end{figure}

\begin{figure}[t]
\centering
\fbox{\parbox{0.95\linewidth}{\centering\vspace{2.5cm}
  \textit{[Placeholder: line chart of cumulative regressed prompts over
  training steps for all four methods.
  X-axis: training step (0--1720).
  Y-axis: cumulative \# of prompts that have shown at least one regression.
  Expected: GRPO grows steeply; \method{} stays near zero throughout.]}
  \vspace{2.5cm}}}
\caption{Cumulative number of training prompts showing at least one
  regression event over training steps.
  \method{}'s Mastery Filter prevents forgetting events from accumulating.}
\label{fig:regression_curve}
\end{figure}

\subsubsection{Hyperparameter Determination}
\label{sec:hparam}

\method{} introduces three hyperparameters beyond the standard GRPO setup:
the EMA decay coefficient $\alpha$, the warmup threshold $K$, and the
mastery threshold $\tau$.
We discuss the design rationale for each and report sensitivity experiments
validating the chosen defaults.

\paragraph{EMA decay coefficient $\alpha$.}
The role of $\alpha$ is to trade off stability against adaptivity:
larger $\alpha$ produces a smoother, more stable baseline that is slower
to reflect recent learning.
The effective sample count $n_{\mathrm{eff}} \approx n/(1-\alpha)$
determines how much of the step-level baseline noise is suppressed.
Table~\ref{tab:alpha_tradeoff} summarizes the properties at each candidate
value.

\begin{table}[h]
\centering
\caption{EMA decay $\alpha$: effective sample count and adaptation speed.
  ``Visits to converge'' = number of visits before $\mu(p)$ tracks the
  true mean to within 10\% of the initial error.}
\label{tab:alpha_tradeoff}
\small
\begin{tabular}{cccp{6cm}}
\toprule
$\alpha$ & $n_{\mathrm{eff}}$ & Visits to converge & Recommended scenario \\
\midrule
0.80 &  40 & $\sim$5 & Large datasets; each prompt appears infrequently per epoch \\
\textbf{0.90} & \textbf{80} & $\sim$10 & \textbf{Default.} Small--medium datasets; each prompt appears $\sim$1$\times$/epoch \\
0.95 & 160 & $\sim$20 & Large datasets with stable prompt distributions \\
0.99 & 800 & $\sim$100 & Very large datasets; very slow adaptation, insensitive to recent regression \\
\bottomrule
\end{tabular}
\end{table}

The choice $\alpha=0.9$ is motivated by the dataset structure.
In our setting with $N=400$ training prompts and batch size 32, each prompt
appears roughly once every $\lceil N/32 \rceil \approx 13$ steps.
An effective window of 10 steps therefore corresponds to approximately
10 consecutive epochs, giving the EMA baseline time to average out
within-epoch randomness while still responding to genuine policy improvement
within a few epochs.
At $\alpha=0.99$, the EMA window spans $\sim$100 steps ($\approx$100 epochs
in this setting), making the baseline almost constant and no better than a
fixed prior.
At $\alpha=0.80$, the window is only 5 steps, barely longer than a single
batch, and the variance reduction over standard GRPO is limited
($5\times$ instead of $10\times$).

To empirically validate $\alpha=0.9$, we run \method{} on
MATH-500$_{\text{400}}$ with Qwen3-1.7B across four $\alpha$ values,
fixing $K=3$.
Table~\ref{tab:hparam_alpha} reports best-checkpoint MATH-500 and GSM8K
pass@1.

\begin{table}[h]
\centering
\caption{Sensitivity to EMA decay $\alpha$ ($K=3$ fixed, Qwen3-1.7B,
  MATH-500$_{\text{400}}$, best checkpoint).}
\label{tab:hparam_alpha}
\small
\begin{tabular}{ccccc}
\toprule
$\alpha$ & $n_{\mathrm{eff}}$ & MATH-500 p@1 & MATH-500 p@8 & \# Regressed prompts \\
\midrule
0.80 &  40 & -- & -- & -- \\
\textbf{0.90} & \textbf{80} & \textbf{0.689} & \textbf{0.800} & \textbf{--} \\
0.95 & 160 & -- & -- & -- \\
0.99 & 800 & -- & -- & -- \\
\bottomrule
\end{tabular}
\end{table}

\paragraph{Warmup threshold $K$.}
Before the EMA baseline has accumulated sufficient observations, it is
initialized from the first batch mean (cold start) and is therefore
subject to single-batch noise.
The warmup threshold $K$ specifies how many visits are required before
the EMA mean is trusted as a baseline; below $K$, the method falls back
to the standard GRPO batch mean.

Table~\ref{tab:warmup_tradeoff} summarizes the implications of different $K$
values.

\begin{table}[h]
\centering
\caption{Warmup threshold $K$: behavior and recommendation.}
\label{tab:warmup_tradeoff}
\small
\begin{tabular}{ccp{7.5cm}}
\toprule
$K$ & Warmup epochs & Effect \\
\midrule
1 & 0 & No warmup; EMA activates immediately on second visit.
        Cold-start mean from one batch ($n=8$ samples) is high-variance;
        early EMA baseline is unreliable. \emph{Not recommended.} \\
\textbf{3} & $\sim$3 & \textbf{Default.} EMA based on 3 visits ($\sim$24 rollouts) before activation;
        mean estimate standard error $\approx 0.177/\sqrt{3} \approx 0.10$.
        Balance between early activation and reliability. \\
5 & $\sim$5 & More conservative; reliable EMA from 5 visits ($\sim$40 rollouts).
        Foregoes improvement for the first 5 epochs.
        Marginal benefit over $K=3$ in practice. \\
10 & $\sim$10 & Very conservative; for the first 10 epochs the method is
        equivalent to standard GRPO. Unnecessarily delays forgetting protection. \\
\bottomrule
\end{tabular}
\end{table}

The choice $K=3$ reflects a compromise: after 3 visits the per-prompt
mean estimate has standard error $\approx \sqrt{\pi(1-\pi)/(3n)} \approx 0.10$
(at $\pi=0.5$, $n=8$), which is roughly half the single-batch error.
This is sufficient for the EMA to provide a more stable baseline than the
batch mean without delaying protection for too long.
Crucially, the Mastery Filter is active from step 1 regardless of $K$---since
it operates on the current-batch mean rather than the EMA---so forgetting
protection is never entirely disabled during warmup.

Table~\ref{tab:hparam_K} reports empirical sensitivity to $K$ on
MATH-500$_{\text{400}}$ with $\alpha=0.9$ fixed.

\begin{table}[h]
\centering
\caption{Sensitivity to warmup threshold $K$ ($\alpha=0.9$ fixed,
  Qwen3-1.7B, MATH-500$_{\text{400}}$, best checkpoint).}
\label{tab:hparam_K}
\small
\begin{tabular}{cccc}
\toprule
$K$ & MATH-500 p@1 & MATH-500 p@8 & \# Regressed prompts \\
\midrule
1  & -- & -- & -- \\
\textbf{3}  & \textbf{0.689} & \textbf{0.800} & \textbf{--} \\
5  & -- & -- & -- \\
10 & -- & -- & -- \\
\bottomrule
\end{tabular}
\end{table}

\paragraph{Mastery threshold $\tau$.}
We fix $\tau=1.0$ (all $n$ rollouts correct) for two reasons.
First, requiring strict unanimity avoids silencing prompts where the model
is not yet reliably correct: a batch with $n-1$ correct and 1 incorrect
rollout still contains a valid contrastive signal---the failed rollout has
advantage $r^{\text{fail}} - \mu_{t-1}(p) < 0$ and the successful rollouts
have positive advantages, providing a well-conditioned update direction.
Setting $\tau < 1.0$ (\eg, $\tau = 7/8 = 0.875$) would suppress this
signal, treating a prompt as mastered when a fraction of rollouts remain
incorrect.
Second, for binary rewards, $\hat{\mu}_t^{(p)} = 1.0$ is the only
configuration where $\hat{\sigma}_t^{(p)} = 0$ exactly, which means
$\tau=1.0$ is precisely the threshold at which std-normalized advantages
would be numerically undefined---making the Filter a necessary and
sufficient condition for numerical safety (Section~\ref{sec:ema_norm}).
We therefore do not treat $\tau$ as a tunable hyperparameter and do not
report a sensitivity sweep for it.

\subsubsection{Degradation-Aware Resampling vs.\ Mastery Filter}
\label{sec:resampling_ablation}

The Mastery Filter provides an implicit difficulty-aware curriculum:
mastered prompts contribute zero gradient, so all learning signal comes
from non-mastered prompts automatically.
Degradation-Aware Resampling (Section~\ref{sec:resampling}) pursues the same
objective in sampling space and is therefore conceptually redundant with
the Filter.
Table~\ref{tab:ablation_resampling} tests whether explicit resampling at
varying $\gamma$ provides any additional gain over the Mastery Filter alone
($\gamma=0$) on MATH-Train$_{\text{2800}}$.
We expect the answer to be negative: the implicit gradient curriculum
is sufficient, and explicit oversampling of hard prompts only increases
overfitting risk.

\begin{table}[t]
\centering
\caption{Degradation-Aware Resampling: effect of $\gamma$ across all benchmarks
  (Qwen2.5-Math-1.5B, MATH-Train$_{\text{2800}}$, best checkpoint).
  $\gamma=0$ is \method{} without resampling.}
\label{tab:ablation_resampling}
\resizebox{\linewidth}{!}{%
\begin{tabular}{ccccccccc}
\toprule
& \multicolumn{2}{c}{\textbf{AIME 2024}} & \multicolumn{2}{c}{\textbf{AIME 2025}}
& \multicolumn{2}{c}{\textbf{MATH-500}} & \multicolumn{2}{c}{\textbf{GSM8K}} \\
\cmidrule(lr){2-3}\cmidrule(lr){4-5}\cmidrule(lr){6-7}\cmidrule(lr){8-9}
$\gamma$ & p@1 & p@8 & p@1 & p@8 & p@1 & p@8 & p@1 & p@8 \\
\midrule
0.0 (no resampling) & -- & -- & -- & -- & -- & -- & -- & -- \\
0.5 & -- & -- & -- & -- & -- & -- & -- & -- \\
1.0 & -- & -- & -- & -- & -- & -- & -- & -- \\
2.0 & -- & -- & -- & -- & -- & -- & -- & -- \\
\bottomrule
\end{tabular}}
\end{table}

\subsubsection{Scale Generalization}
\label{sec:scale}

Table~\ref{tab:scale} tests whether the gains of \method{} over GRPO
generalize across training set sizes.
Qwen3-1.7B results on the two smaller splits are populated from completed
experiments; Qwen2.5-Math-1.5B results on MATH-Train$_{\text{2800}}$
are to be filled once all checkpoints are evaluated.

\begin{table}[t]
\centering
\caption{Scale generalization: \method{} vs.\ GRPO across training set sizes
  and model configurations (best checkpoint, $\Delta = \text{\method{}} - \text{GRPO}$).}
\label{tab:scale}
\resizebox{\linewidth}{!}{%
\begin{tabular}{llcccccccc}
\toprule
& & \multicolumn{2}{c}{\textbf{AIME 2024}} & \multicolumn{2}{c}{\textbf{AIME 2025}}
& \multicolumn{2}{c}{\textbf{MATH-500}} & \multicolumn{2}{c}{\textbf{GSM8K}} \\
\cmidrule(lr){3-4}\cmidrule(lr){5-6}\cmidrule(lr){7-8}\cmidrule(lr){9-10}
\textbf{Dataset ($N$)} & \textbf{Model}
  & $\Delta$p@1 & $\Delta$p@8 & $\Delta$p@1 & $\Delta$p@8
  & $\Delta$p@1 & $\Delta$p@8 & $\Delta$p@1 & $\Delta$p@8 \\
\midrule
MATH-500$_{\text{400}}$ (400)
  & Qwen3-1.7B           & -- & -- & -- & -- & +0.019 & +0.010 & -- & -- \\
GSM8K$_{\text{1000}}$ (1{,}000)
  & Qwen3-1.7B           & -- & -- & -- & -- & -- & -- & +0.007 & +0.015 \\
\midrule
\multirow{3}{*}{MATH-Train$_{\text{2800}}$ (2{,}800)}
  & Qwen3-1.7B           & -- & -- & -- & -- & -- & -- & -- & -- \\
  & Qwen2.5-Math-1.5B    & -- & -- & -- & -- & -- & -- & -- & -- \\
  & Qwen3-8B           & -- & -- & -- & -- & -- & -- & -- & -- \\
  & Qwen2.5-Math-7B-Base & -- & -- & -- & -- & -- & -- & -- & -- \\
\bottomrule
\end{tabular}}
\end{table}

% ──────────────────────────────────────────────────────────────────────────────
\subsection{Training Dynamics Analysis}
\label{sec:dynamics}
% ──────────────────────────────────────────────────────────────────────────────

Beyond benchmark numbers, we analyze the internal statistics of \method{}
to understand why it mitigates forgetting.

\paragraph{EMA mean distribution evolution.}
Figure~\ref{fig:ema_dist} visualizes how the distribution of per-prompt
EMA means $\{\mu_t(p)\}_{p \in \mathcal{D}}$ evolves over training.
We expect a rightward shift (overall improvement) accompanied by emerging
bimodality: easy prompts converge toward $\mu=1.0$ while hard prompts
remain concentrated near lower values.

\begin{figure}[t]
\centering
\fbox{\parbox{0.95\linewidth}{\centering\vspace{3cm}
  \textit{[Placeholder: histogram or violin plot of per-prompt EMA mean
  $\mu_t(p)$ at training steps $t \in \{0, 400, 800, 1200, 1720\}$.
  X-axis: EMA mean value (0--1). Y-axis: density.
  Color-coded curves for each checkpoint.
  Expected: distribution shifts right and becomes bimodal over time,
  separating mastered ($\mu\to 1$) from persistently hard prompts.]}
  \vspace{3cm}}}
\caption{Distribution of per-prompt EMA means $\mu_t(p)$ at five training
  checkpoints. The rightward shift indicates overall learning; the bimodal
  structure confirms that the bank correctly differentiates mastered from
  hard prompts.}
\label{fig:ema_dist}
\end{figure}

\paragraph{Mastery Filter trigger rate.}
Figure~\ref{fig:filter_rate} tracks the fraction of prompts per step for
which the Mastery Filter fires ($\hat{\mu}_t^{(p)} = 1.0$, all rollouts
correct).
A monotonically increasing rate is the expected signal: as the model learns,
more prompts become fully mastered and their gradients are suppressed.

\begin{figure}[t]
\centering
\fbox{\parbox{0.95\linewidth}{\centering\vspace{2.5cm}
  \textit{[Placeholder: line chart of Mastery Filter trigger rate
  (fraction of batch prompts with $\hat\mu_t^{(p)}=1.0$) over training steps.
  X-axis: training step (0--1720). Y-axis: trigger rate (0--1).
  Shaded region: rolling-window standard deviation across steps.
  Expected: monotonic increase from near 0 to 0.3--0.5 by epoch~20.]}
  \vspace{2.5cm}}}
\caption{Mastery Filter trigger rate over training. The monotonically rising
  rate confirms that the filter increasingly suppresses gradients on mastered
  prompts, protecting acquired knowledge throughout training.}
\label{fig:filter_rate}
\end{figure}

\paragraph{Advantage magnitude distribution.}
Figure~\ref{fig:adv_magnitude} compares the distribution of non-zero advantage
magnitudes $|A_i^{(p)}|$ between GRPO and \method{} over training.
Because GRPO normalizes by batch std, its advantages are tightly bounded.
\method{}'s unnormalized advantages $r - \mu$ preserve natural scale:
near-mastered prompts emit small advantages while hard prompts emit larger
ones, providing an implicit difficulty-adaptive weighting without explicit
resampling.

\begin{figure}[t]
\centering
\fbox{\parbox{0.95\linewidth}{\centering\vspace{2.5cm}
  \textit{[Placeholder: violin or box plots comparing $|A_i^{(p)}|$ distributions
  for GRPO vs.\ \method{} at training checkpoints step $\in \{400, 800, 1200, 1720\}$.
  X-axis: checkpoint. Y-axis: $|A_i^{(p)}|$.
  Expected: GRPO advantages cluster uniformly near $[0.1, 0.5]$;
  \method{} advantages are smaller for high-$\mu$ prompts and larger
  for low-$\mu$ prompts, showing difficulty-adaptive concentration.]}
  \vspace{2.5cm}}}
\caption{Distribution of advantage magnitudes $|A_i^{(p)}|$ for GRPO vs.\
  \method{} across training checkpoints.
  \method{}'s unnormalized formulation naturally concentrates gradient
  energy on harder, less-mastered prompts.}
\label{fig:adv_magnitude}
\end{figure}

% ──────────────────────────────────────────────────────────────────────────────
\subsection{Multi-Model Generalization}
\label{sec:largescale}
% ──────────────────────────────────────────────────────────────────────────────

To test whether \method{} generalizes beyond the primary 1.5B model, we
evaluate on Qwen3-1.7B and Qwen2.5-Math-7B-Base under the same
MATH-Train$_{\text{2800}}$ setup.
A stronger base model raises overall accuracy, making it a meaningful test of
whether forgetting suppression remains beneficial when the model already
solves a larger fraction of training prompts.
Results follow the same evaluation protocol as Section~\ref{sec:main_results}.

\begin{table}[t]
\centering
\caption{Results with \textbf{Qwen3-1.7B} on MATH-Train$_{\text{2800}}$,
  \textbf{final} checkpoint (step~1720).
  Best per column in \textbf{bold}.}
\label{tab:qwen3_1.7b_results}
\resizebox{\linewidth}{!}{%
\begin{tabular}{lcccccccccccc}
\toprule
& \multicolumn{3}{c}{\textbf{AIME 2024}} & \multicolumn{3}{c}{\textbf{AIME 2025}}
& \multicolumn{3}{c}{\textbf{MATH-500}} & \multicolumn{3}{c}{\textbf{GSM8K}} \\
\cmidrule(lr){2-4}\cmidrule(lr){5-7}\cmidrule(lr){8-10}\cmidrule(lr){11-13}
\textbf{Method} & p@1 & p@4 & p@8 & p@1 & p@4 & p@8 & p@1 & p@4 & p@8 & p@1 & p@4 & p@8 \\
\midrule
GRPO             & -- & -- & -- & -- & -- & -- & -- & -- & -- & -- & -- & -- \\
Dr.~GRPO         & -- & -- & -- & -- & -- & -- & -- & -- & -- & -- & -- & -- \\
DAPO             & -- & -- & -- & -- & -- & -- & -- & -- & -- & -- & -- & -- \\
\method{} (ours) & -- & -- & -- & -- & -- & -- & -- & -- & -- & -- & -- & -- \\
\bottomrule
\end{tabular}}
\end{table}

\begin{table}[t]
\centering
\caption{Results with \textbf{Qwen2.5-Math-7B-Base} on MATH-Train$_{\text{2800}}$,
  \textbf{final} checkpoint.
  Best per column in \textbf{bold}.}
\label{tab:7b_results}
\resizebox{\linewidth}{!}{%
\begin{tabular}{lcccccccccccc}
\toprule
& \multicolumn{3}{c}{\textbf{AIME 2024}} & \multicolumn{3}{c}{\textbf{AIME 2025}}
& \multicolumn{3}{c}{\textbf{MATH-500}} & \multicolumn{3}{c}{\textbf{GSM8K}} \\
\cmidrule(lr){2-4}\cmidrule(lr){5-7}\cmidrule(lr){8-10}\cmidrule(lr){11-13}
\textbf{Method} & p@1 & p@4 & p@8 & p@1 & p@4 & p@8 & p@1 & p@4 & p@8 & p@1 & p@4 & p@8 \\
\midrule
GRPO             & -- & -- & -- & -- & -- & -- & -- & -- & -- & -- & -- & -- \\
Dr.~GRPO         & -- & -- & -- & -- & -- & -- & -- & -- & -- & -- & -- & -- \\
DAPO             & -- & -- & -- & -- & -- & -- & -- & -- & -- & -- & -- & -- \\
\method{} (ours) & -- & -- & -- & -- & -- & -- & -- & -- & -- & -- & -- & -- \\
\bottomrule
\end{tabular}}
\end{table}

% ══════════════════════════════════════════════════════════════════════════════
% Appendix
% ══════════════════════════════════════════════════════════════════════════════
\appendix

\section{Implementation Details}
\label{app:impl}

\paragraph{Full hyperparameter table.}
Table~\ref{tab:hyperparams_full} lists all hyperparameters used across
experiments.

\begin{table}[h]
\centering
\caption{Full hyperparameter listing used in all experiments.}
\label{tab:hyperparams_full}
\small
\begin{tabular}{ll}
\toprule
\textbf{Hyperparameter} & \textbf{Value} \\
\midrule
Rollout count $n$                          & 8 \\
LoRA rank / $\alpha_{\mathrm{LoRA}}$       & 64 / 32 \\
Optimizer                                  & AdamW, lr $= 10^{-5}$ \\
Batch size                                 & 32 \\
Training epochs                            & 20 \\
Max response length                        & 8{,}192 tokens \\
KL coefficient $\beta$                     & 0.001 (low-variance KL estimator) \\
EMA decay $\alpha$                         & 0.9 \\
Warmup threshold $K$                       & 3 \\
Mastery threshold $\tau$                   & 1.0 \\
Eval temperature                           & 0.6 \\
Eval samples per problem                   & 8 \\
\bottomrule
\end{tabular}
\end{table}

\paragraph{PromptMemoryBank complexity and integration.}
The PromptMemoryBank stores two floating-point scalars per prompt (EMA mean
and visit count), requiring $O(N)$ additional memory—negligible relative to
model parameters.
It is serialized alongside model checkpoints to ensure continuity across
training resumptions.
Integration into an existing GRPO framework requires three modifications:
(1) initializing the memory bank at training start,
(2) intercepting the advantage computation step to apply Eq.~\eqref{eq:final_adv},
and (3) updating the bank after each rollout batch.
No changes to the model architecture, optimizer, or rollout infrastructure
are required.

% ─────────────────────────────────────────────────────────────────────────────

\section{Degradation-Aware Resampling}
\label{app:resampling}

We explored an explicit resampling scheme that upweights hard and regressing
prompts in the sampling distribution.
As argued in Section~\ref{sec:mastery}, the Mastery Filter already achieves
this goal implicitly: mastered prompts contribute zero gradient, so all
learning signal originates from non-mastered prompts automatically.
Explicit resampling is therefore conceptually redundant with the Filter,
and our ablations (Section~\ref{sec:resampling_ablation}) confirm it
provides no additional benefit and is actively harmful on small datasets.
We describe it here for completeness.

Define the \emph{difficulty score} $d_t(p) = \max(0, 1 - \mu_t(p))$
and the \emph{degradation score} $g_t(p) = \max(0, \mu_0(p) - \mu_t(p))$,
both zero for $c(p) < K$.
The sampling weight for prompt $p$ at step $t$ is
\begin{equation}
  w_t(p) = 1 + \gamma\!\left[d_t(p) + g_t(p)\right],
  \label{eq:weight}
\end{equation}
where $\gamma=0$ recovers uniform sampling.
The difficulty score targets prompts the model has never reliably solved;
the degradation score targets prompts that were previously mastered but
have since regressed—both desirable properties for forgetting mitigation.

\paragraph{Why it fails on small datasets.}
With $N=400$ training prompts, a hard subset of $\sim$50 prompts receives
$2\times$--$3\times$ more samples under $\gamma=1$.
This effectively reduces the diversity of the training distribution by a
factor of 2--3, causing the policy to overfit rapidly to the hard subset
while losing generalization on the rest.
The implicit curriculum of the Mastery Filter (gradient suppression on
mastered prompts) achieves difficulty concentration in \emph{gradient space}
without altering the sampling distribution, avoiding this overfitting failure.
We recommend Degradation-Aware Resampling only when $N > 1{,}000$.

\end{document}
